\section{Un protocole synthétique figé et traçable}
\label{sec:protocol}

\subsection{Deux images aux structures complémentaires}

Les images sont construites par \path{datasets.py}, sans téléchargement.
La première, identifiée par \path{geometric_shapes}, comporte un fond d'intensité
$0{,}15$, une moitié droite à $0{,}32$, un rectangle à $0{,}72$ et un
disque à $0{,}92$. Avec $H=W=128$, le rectangle utilise les lignes
d'indices $H//6$ à $H//2-1$ et les colonnes $W//10$ à $2W//5-1$.
Le disque a pour centre $(2H/3,3W/4)$ en coordonnées ligne/colonne et
pour rayon $\min(H,W)/6$. Cette image teste principalement des plateaux
et des ruptures nettes, horizontales, verticales et courbes.

La seconde, identifiée par \path{ramps_and_edges}, introduit une variation régulière.
En coordonnées normalisées $x=j/(W-1)$ et $y=i/(H-1)$, son fond est
$0{,}08+0{,}48x$ ; un saut de $0{,}16$ est ajouté pour $y\geq0{,}55$.
Une ellipse centrée en $(0{,}28,0{,}70)$, de demi-axes $0{,}16$ et
$0{,}18$, est remplacée par l'intensité $0{,}88$. Enfin, la bande
$|y-(0{,}12+0{,}68x)|\leq0{,}025$ est fixée à $0{,}04$.
Le code borne cette \emph{construction de vérité terrain} dans $[0,1]$.
Cela ne doit pas être confondu avec une troncature du bruit ou des sorties,
qui n'est jamais pratiquée. L'image confronte le schéma directionnel à
une rampe, une ellipse et une structure oblique ; elle ne constitue pas
une validation exhaustive de l'invariance par rotation.

\subsection{Grille de paramètres}

Le plan étendu a été construit \emph{après} l'exploration rapide de
phase~3, puis enregistré avant le calcul de phase~4. Ce n'est pas une
évaluation prospective sur un ensemble d'images indépendant. Les
paramètres de \path{configs/extended.json} sont :
\begin{table}[htbp]
 \centering
 \caption{Protocole étendu figé avant exécution ; temps de diffusion en
 convention pixel, sans unité physique en secondes.}
 \label{tab:protocol}
 \begin{tabular}{ll}
  \toprule
  Élément & Valeurs \\
  \midrule
  Images et taille & Deux images synthétiques, $128\times128$ \\
  Bruit gaussien & $\sigma\in\{0{,}025;\,0{,}05;\,0{,}10\}$ \\
  Graines & $0,1,\ldots,9$ \\
  Pas spatiaux et temporel & $\dx=\dy=1$, $\dt=0{,}20$ \\
  Checkpoints $n$ & $0,1,2,5,10,20,40,80$ \\
  Conductions PM & Exponentielle et rationnelle \\
  Valeurs de $K$ & $0{,}025;\,0{,}05;\,0{,}10;\,0{,}20$ \\
  Plage des métriques & $L=1$ \\
  \bottomrule
 \end{tabular}
\end{table}

Pour chaque tuple image/$\sigma$/graine, une unique observation
$f=u_{\mathrm{true}}+\sigma\xi$ est générée, avec des composantes
gaussiennes standard indépendantes à l'intérieur de $\xi$. Ce même
tableau est fourni à la chaleur et aux huit variantes PM. L'observation
non traitée sert de référence commune à $n=0$.

Le générateur est réinitialisé avec la même graine pour chaque tuple.
Comme toutes les images ont la même taille, une graine donnée réutilise
le même champ gaussien standardisé entre images et niveaux de bruit.
Les comparaisons entre méthodes sont donc appariées de façon contrôlée,
mais les six cas image/$\sigma$ ne sont pas entièrement indépendants.
Les dix graines représentent dix réalisations de bruit sur chaque image,
non dix images différentes.

\subsection{Comptages et budgets comparables}

Le plan contient $2\times3\times10=60$ observations et
$60\times(1+2\times4)=540$ trajectoires. Chaque trajectoire est poursuivie
jusqu'à $n=80$, avec copies aux huit checkpoints ; il ne s'agit pas de
relancer un solveur pour chaque checkpoint. Une observation contribue
une ligne bruitée et $9\times8$ lignes de solveur, soit $73$ lignes.
Les fichiers contiennent ainsi $60\times73=4\,380$ lignes de métriques
et $6\times73=438$ groupes, chacun sur dix graines.

Les tableaux à checkpoint commun comparent le même nombre d'itérations
et le même $t=n\dt$. Cette égalité de budget numérique ne signifie pas
que les conductances produisent le même lissage effectif ni la même durée
machine. Tous les paramètres prédéfinis restent disponibles, y compris
les configurations défavorables. Les tables exhaustives sont séparées
des tableaux ciblés du corps du rapport.

\subsection{Cas représentatif et calcul pilote}

Avant exécution, le cas visuel a été fixé à
\path{geometric_shapes}, $\sigma=0{,}10$, graine~$0$, $n=20$ et
$K=0{,}10$ pour les deux PM. Ce choix n'a pas été modifié pour montrer
a posteriori le meilleur résultat. Ses figures décrivent une réalisation
individuelle, contrairement aux courbes agrégées sur dix graines.

Un pilote séparé, portant sur cette image, ce bruit et cette graine, mais
avec toute la grille de méthodes, $K$ et checkpoints, a pris
$6{,}8015153$~s. L'estimation conservative $60$ fois cette durée était
$408{,}090918$~s, inférieure au seuil préalable de $1\,200$~s. Le plan
complet n'a donc pas été réduit. Le pilote n'est pas ajouté comme une
onzième réalisation aux statistiques de l'étude.
